\problemname{JuQueen}
\emph{JuQueen} is the super computer with the best performance allover Germany.
It is on rank $8$ in the famous top$500$ list with its $458\,752$ cores.
It draws a lot of energy (up to $2\,301$~kW), so we want to reduce that by underclocking the unused cores.

The cluster scheduling algorithm which is in charge of distributing jobs over the nodes and cores of a cluster will issue the following speedstepping commands:
\begin{itemize}
	\item \texttt{change X S} changes the frequency of core \texttt{X} by \texttt{S} steps
	\item \texttt{groupchange A B S} changes the frequency of every core in range \texttt{[A,B]} by \texttt{S} steps
	\item \texttt{state X} returns the current state of core \texttt{X}
\end{itemize}
To be safe for the future, your program should be able to handle $4\,587\,520$ cores.
The initial frequency for each core is $0$.

\begin{Input}
    The input contains a single test case.
    It starts with a line containing three integers $C$, $N$, and $O$,
    where $C$ is the number of cores ($1 \leq C \leq 4\,587\,520$) to manage, $N$ is the number of frequency steps for each core ($1 \leq N \leq 10\,000$) and $O$ is the number of operations in the test program ($1 \leq O \leq 50\,000$).
    Then follow $O$ lines, each containing one command as described above.
    $X$, $A$ and $B$ are $0$-based IDs of the cores ($0 \leq A, B, X < C$; $A \leq B$).
    $S$ is an integer number of steps, possibly negative ($-N \leq S \leq +N$).
    Both, the \texttt{change} and the \texttt{groupchange} command will increase (or decrease) in single steps and stop as soon as \textbf{one} core in the group reaches the minimal ($0$) or maximal frequency ($N$).
\end{Input}

\begin{Output}
    Output one line for every operation in the input.
    For \texttt{change} and \texttt{groupchange} print the changed
    number of steps, for \texttt{state} print the current state.
\end{Output}
